\documentclass[12pt,a4paper]{article}
\usepackage{amsmath,amsfonts,amssymb}
\usepackage{graphicx}
\usepackage{hyperref}
\usepackage{enumitem}
\usepackage{geometry}
\geometry{margin=2.5cm}
\emergencystretch=2em

\begin{document}

\begin{center}
    {\Large \textbf{Master Thesis Project in Quantum Information Science and Technology}}\\[0.3cm]
    {\large \textbf{Many-Body Entanglement and Quantum Fisher Information of up to Four Trapped Electrons: Multi-Configuration Time-Dependent Hartree--Fock and Reinforcement Learning}}\\[0.2cm]
    \vspace{0.3cm}
    {Duration: One Year (60 ECTS)}
\end{center}

\vspace{1cm}

\section*{1. Introduction}

Electrons trapped above the surface of superfluid helium are a
candidate platform for quantum computing and quantum sensing~[1--3].
Each electron sits in its own electrostatic well, the two lowest
states of its motion define a qubit, and the unscreened Coulomb
interaction between electrons in different wells couples them.  The
interaction entangles the electrons, and this entanglement is the
resource behind two-qubit gates and behind sensors that reach a
precision beyond that of independent particles.

With two electrons, entanglement is fully described by one number,
for example the von Neumann entropy of one electron, and this is the
quantity studied in our earlier work~[2].  With three and four
electrons this is no longer the case.  There are several inequivalent
kinds of entanglement, such as those of the GHZ, W and cluster states,
which cannot be transformed into each other by local operations~[4,5],
and no single measure distinguishes them.  At the same time the
quantum Fisher information (QFI), which sets the ultimate precision of
a sensor, is itself a witness of how many particles are entangled with
each other~[6--8].  Understanding which kind of many-body entanglement
a device can produce, and which kind is useful for a given task,
requires several measures side by side.

In the PhD work of Oskar Leinonen~[3,12], the multi-configuration
time-dependent Hartree and Hartree--Fock method~[9--11], MCTDH(F), has
been developed for registers of up to four wells.  It propagates
the full correlated wave function of all electrons, including the
states outside the qubit subspace, while the well parameters and
external fields change in time.  This thesis project uses that
framework as its solver and adds two elements: a systematic analysis
of many-body entanglement with measures beyond the von Neumann
entropy, and reinforcement learning (RL)~[13--16] to find static
configurations and time-dependent control fields that maximize the
QFI and selected entanglement measures.  The coupling to an
environment and decoherence are not part of this project: the system
is closed and the focus is on the entanglement itself.

\section*{2. Physical System and Scientific Motivation}

The system is the register of anharmonic wells of Refs.~[3,12] with
up to four electrons, one per well.  With $\hbar=m_e=1$,
\[
  \hat H(t)=\sum_{\sigma=1}^{N}\Bigl[-\frac12\frac{\partial^2}{\partial x_\sigma^2}
  +V^{\sigma}(x_\sigma;t)+\mathcal E_\sigma(t)\,x_\sigma\Bigr]
  +\sum_{\sigma<\tau}\frac{\kappa}{|x_\sigma-x_\tau|},
\]
\[
  V^{\sigma}(x)=\tfrac12(\omega^{\sigma})^2(x-x_c^{\sigma})^2
  +\tfrac13\beta^{\sigma}(\omega^{\sigma})^2(x-x_c^{\sigma})^4 ,
\]
where the frequencies $\omega^\sigma$ and anharmonicities
$\beta^\sigma$ are set by the gate electrodes and $\mathcal E_\sigma$
is an electric field acting on well $\sigma$.  The controls are of two
kinds.
\begin{itemize}
  \item \textbf{Time-independent fields.}  Static values of
  $\omega^\sigma$, $\beta^\sigma$ and of a static electric field.
  They define the operating points of the device and the entanglement
  of its eigenstates, which is largest near avoided crossings.
  \item \textbf{Time-dependent fields.}  Ramps of the well parameters
  between operating points, which realize two-qubit gates such as
  $\sqrt{i\mathrm{SWAP}}$~[3,12], and resonant microwave fields
  $\mathcal E_\sigma(t)$ on single wells, which realize single-qubit
  rotations.  Sequences of these generate entanglement between
  several electrons.
\end{itemize}

In MCTDH(F) the wave function is expanded in products of
time-dependent orbitals of the individual wells,
\[
  |\Psi(t)\rangle=\sum_{i_1\dots i_N}C_{i_1\dots i_N}(t)\,
  |\varphi^{1}_{i_1}(t)\rangle\otimes\cdots\otimes|\varphi^{N}_{i_N}(t)\rangle .
\]
With one electron per well the wells are distinguishable subsystems,
and the reduced density matrix $\rho^\sigma$ of well $\sigma$ is the
one-body density matrix that the method computes at every time step.
Reduced density matrices of groups of wells follow from the
coefficients $C$.  The entanglement measures below are therefore
available at little extra cost, for the full wave function and not
only for its projection on the qubit states.
\begin{itemize}
  \item \textbf{Bipartite measures.}  For every division of the wells
  into two groups $A|B$: the von Neumann entropy, the R\'enyi
  entropies $S_\alpha=(1-\alpha)^{-1}\ln\mathrm{Tr}\,\rho_A^\alpha$
  and the entanglement spectrum~[4,17].
  \item \textbf{Entanglement of two wells inside a larger system.}
  The state of two wells is then mixed, and the entropy is no longer
  a measure of entanglement.  The negativity
  $\mathcal N=(\|\rho_{\sigma\tau}^{T_\sigma}\|_1-1)/2$~[18] and, for
  qubits, the concurrence~[19] are, and they obey monogamy
  relations that define the three-tangle~[20].
  \item \textbf{Multipartite measures.}  The global entanglement of
  Meyer and Wallach,
  $Q=2\bigl[1-\tfrac1N\sum_\sigma\mathrm{Tr}(\rho^\sigma)^2\bigr]$~[21],
  the genuine multipartite concurrence
  $C_{\rm GME}=\min_{A|B}\sqrt{2\,[1-\mathrm{Tr}\,\rho_A^2]}$~[22],
  which vanishes unless all electrons are entangled, the geometric
  measure of entanglement~[23], and the classification of four-qubit
  states~[5].
  \item \textbf{Quantum Fisher information.}  For a pure state and a
  parameter $\lambda$ of the Hamiltonian,
  $F_Q=4\bigl[\langle\partial_\lambda\Psi|\partial_\lambda\Psi\rangle
  -|\langle\Psi|\partial_\lambda\Psi\rangle|^2\bigr]$, which reduces
  to four times the variance of the generator $\hat G$ when the
  parameter is imprinted as $e^{-i\lambda\hat G}$~[6].  For $N$ qubits
  and $\hat G=\tfrac12\sum_\sigma\hat Z_\sigma$, separable states give
  $F_Q\le N$, a value $F_Q>kN$ proves that at least $k+1$ particles
  are entangled, and the maximum $N^2$, here 16, is reached by the GHZ
  state~[7,8,24].
\end{itemize}

The project is organized around three questions:
\begin{enumerate}
  \item Which kinds of many-body entanglement do the eigenstates and
  the gate dynamics of two, three and four interacting electrons
  produce, and what do the measures beyond the von Neumann entropy
  reveal that the entropy alone does not?  How much do the answers
  change when the states outside the qubit subspace are kept?
  \item Can an RL agent find static operating points and
  time-dependent control fields that maximize the QFI, and how does
  it compare with conventional optimization and with sequences of
  calibrated gates?
  \item How are the QFI and the entanglement measures related for
  the states that the device can reach: do controls that maximize
  the QFI also maximize the genuine multipartite entanglement, and
  which fields can be sensed with an advantage over independent
  electrons?
\end{enumerate}

\section*{3. Objectives of the Thesis}

\subsection*{Objective 1: Theory and reproduction}
\begin{itemize}
  \item Study the MCTDH(F) method, entanglement theory for bipartite
  and multipartite systems, quantum parameter estimation and
  reinforcement learning.
  \item Learn the MCTDH(F) codes of O.~Leinonen and reproduce, for two
  wells, the operating points and the $\sqrt{i\mathrm{SWAP}}$ gate of
  Refs.~[3,12], with a check against an exact grid calculation.
\end{itemize}

\subsection*{Objective 2: Entanglement toolbox}
\begin{itemize}
  \item Implement the measures of Sec.~2 for MCTDH(F) wave functions:
  reduced density matrices of one and two wells, entropies and
  entanglement spectra for all bipartitions, negativity, concurrence
  and tangles, global and genuine multipartite entanglement, the
  geometric measure, and the QFI for static and time-dependent
  parameters.
  \item Evaluate each measure both for the full wave function and
  for its projection on the qubit states, together with the weight
  outside the qubit subspace (leakage).
  \item Validate on states with known values (Bell, GHZ, W, Dicke
  and cluster states).
\end{itemize}

\subsection*{Objective 3: Entanglement of two, three and four electrons}
\begin{itemize}
  \item Time-independent fields: map the entanglement of the
  low-lying eigenstates as a function of the well parameters and of
  a static electric field, in particular near the avoided crossings
  used as target operating points.
  \item Time-dependent fields: follow all measures during single
  gates and during gate sequences that produce GHZ- and W-type states
  of three and four electrons; identify how entanglement spreads
  through the register and where it is lost to leakage.
  \item Compare the measures with each other and establish which of
  them separate the different classes of states.
\end{itemize}

\subsection*{Objective 4: Reinforcement learning with time-independent fields}
\begin{itemize}
  \item Formulate the choice of operating point as an RL environment:
  the action is the vector of well parameters and static fields, the
  reward is the QFI (or a chosen entanglement measure) of the target
  eigenstate, with a penalty for loss of the qubit structure.
  \item Train an actor--critic agent with Proximal Policy
  Optimization~[14] and compare with the derivative-free
  optimization used in Ref.~[12].
\end{itemize}

\subsection*{Objective 5: Reinforcement learning with time-dependent fields}
\begin{itemize}
  \item Extend the environment to a sequential problem: at each step
  the agent chooses the ramp and hold times, the pair of wells to
  couple and the drive amplitudes; the state is propagated with
  MCTDH(F); the reward at the final time is
  \[
     R=\frac{F_Q}{N^2}-w\,(1-s),
  \]
  with $s$ the population that remains in the qubit subspace.
  \item Optimize the QFI for static and for oscillating
  fields~[16,25], for
  two, three and four electrons, and compare the learned protocols
  with gate sequences built from separately calibrated gates and with
  gradient-based optimal control~[16,26].
  \item To keep the cost manageable, train on a hierarchy of models:
  an effective few-level model extracted from the MCTDH(F)
  eigenstates, MCTDH(F) with mean-field spectator wells, and the full
  correlated propagation for validation and final optimization.
  Quantify what is lost when a protocol learned on a reduced model is
  run on the full one.
\end{itemize}

\subsection*{Objective 6: QFI and many-body entanglement}
\begin{itemize}
  \item Characterize the states found by the agent with the full set
  of measures and determine their entanglement depth from the QFI.
  \item Use different measures as reward and compare the resulting
  states: QFI, genuine multipartite concurrence and global
  entanglement do not have the same maximizers.
  \item Study the sensitivity of the optimized protocols to errors
  in the control parameters (closed system).
\end{itemize}

\subsection*{Objective 7: Possible extensions}
\begin{itemize}
  \item Four electrons in fewer wells (two per well), where the
  electrons in a well are indistinguishable, antisymmetry is required
  and mode and particle entanglement must be distinguished~[27].
  \item Three electrons per well as in Ref.~[12].
  \item Measurement protocols that reach the QFI bound (classical
  Fisher information of realistic readouts).
\end{itemize}

\section*{4. Methods and Tools}

\begin{itemize}
  \item \textbf{Many-body dynamics:} the multi-species MCTDH(F) codes
  developed by O.~Leinonen (GPU-based), extended in this project with
  the entanglement and QFI analysis.
  \item \textbf{Reinforcement learning:} PyTorch or JAX with a
  Gymnasium-type environment that wraps the solvers; the PPO
  implementation of Ref.~[28] serves as a reference.
  \item \textbf{Useful background:} quantum mechanics, many-body
  physics (FYS4480), quantum computing (FYS5419) and machine learning
  (FYS-STK4155).
\end{itemize}
All codes and results will be documented in Jupyter notebooks so that
the thesis can be reproduced.

\section*{5. Work Plan and Milestones (One-Year Timeline)}

\subsection*{Semester 1 (Months 1--6)}

\textbf{Months 1--2: Theory and reproduction}
\begin{itemize}
    \item Literature study; MCTDH(F) codes; reproduction of the
    two-well operating points and gate.
\end{itemize}
\textbf{Deliverable:} Short report with the reproduced two-electron
results.

\bigskip

\textbf{Months 3--4: Entanglement toolbox}
\begin{itemize}
    \item Implementation and validation of the measures; first
    application to two and three electrons.
\end{itemize}
\textbf{Deliverable:} Tested analysis module for MCTDH(F) wave
functions.

\bigskip

\textbf{Months 5--6: Entanglement of up to four electrons}
\begin{itemize}
    \item Static maps and gate dynamics for three and four electrons;
    comparison of the measures.
\end{itemize}
\textbf{Deliverable:} Survey of the many-body entanglement produced
by the device.

\subsection*{Semester 2 (Months 7--12)}

\textbf{Months 7--8: RL with time-independent fields}
\begin{itemize}
    \item RL environment and PPO agent for the operating points;
    comparison with derivative-free optimization; effective models
    for fast training.
\end{itemize}
\textbf{Deliverable:} Operating points that maximize the QFI.

\bigskip

\textbf{Months 9--10: RL with time-dependent fields}
\begin{itemize}
    \item Sequential control of ramps and drives; QFI optimization
    for two, three and four electrons; validation with the full
    propagation.
\end{itemize}
\textbf{Deliverable:} Learned control protocols and their comparison
with calibrated gate sequences.

\bigskip

\textbf{Month 11: Analysis and consolidation}
\begin{itemize}
    \item Relation between QFI and the entanglement measures for the
    optimized states; robustness; one extension depending on
    progress.
\end{itemize}

\bigskip

\textbf{Month 12: Writing and final submission}
\begin{itemize}
    \item Completion of the thesis manuscript (methods and theory
    chapters are written continuously during the year) and
    preparation of the final presentation.
\end{itemize}
\textbf{Deliverable:} Completed thesis and defense presentation.

\section*{6. Expected Outcomes}

\begin{itemize}
    \item An entanglement and QFI analysis module for MCTDH(F) wave
    functions, covering bipartite and multipartite measures.
    \item A survey of the many-body entanglement of two, three and
    four interacting electrons under static and time-dependent
    fields, including the role of states outside the qubit subspace.
    \item An RL framework with MCTDH(F) as environment, and control
    protocols that maximize the QFI.
    \item A quantitative account of how the QFI and the different
    entanglement measures are related in a realistic device model.
    \item A master thesis and, depending on the results, a
    contribution to a scientific article.
\end{itemize}

\section*{7. References}

\begin{enumerate}
    \item P.~M.~Platzman and M.~I.~Dykman, \textit{Quantum Computing with Electrons Floating on Liquid Helium}, Science \textbf{284}, 1967 (1999).
    \item N.~R.~Beysengulov \textit{et al.}, \textit{Coulomb Interaction-Driven Entanglement of Electrons on Helium}, PRX Quantum \textbf{5}, 030324 (2024).
    \item O.~Leinonen \textit{et al.}, \textit{Design and dynamics of two-qubit gates with motional states of electrons on helium}, Phys.\ Rev.\ A \textbf{113}, 032624 (2026).
    \item R.~Horodecki, P.~Horodecki, M.~Horodecki and K.~Horodecki, \textit{Quantum entanglement}, Rev.\ Mod.\ Phys.\ \textbf{81}, 865 (2009).
    \item F.~Verstraete, J.~Dehaene, B.~De~Moor and H.~Verschelde, \textit{Four qubits can be entangled in nine different ways}, Phys.\ Rev.\ A \textbf{65}, 052112 (2002).
    \item S.~L.~Braunstein and C.~M.~Caves, \textit{Statistical distance and the geometry of quantum states}, Phys.\ Rev.\ Lett.\ \textbf{72}, 3439 (1994); V.~Giovannetti, S.~Lloyd and L.~Maccone, \textit{Quantum metrology}, Phys.\ Rev.\ Lett.\ \textbf{96}, 010401 (2006).
    \item L.~Pezz\'e and A.~Smerzi, \textit{Entanglement, Nonlinear Dynamics, and the Heisenberg Limit}, Phys.\ Rev.\ Lett.\ \textbf{102}, 100401 (2009).
    \item G.~T\'oth, \textit{Multipartite entanglement and high-precision metrology}, Phys.\ Rev.\ A \textbf{85}, 022322 (2012); P.~Hyllus \textit{et al.}, \textit{Fisher information and multiparticle entanglement}, Phys.\ Rev.\ A \textbf{85}, 022321 (2012).
    \item H.-D.~Meyer, U.~Manthe and L.~S.~Cederbaum, \textit{The multi-configurational time-dependent Hartree approach}, Chem.\ Phys.\ Lett.\ \textbf{165}, 73 (1990); M.~H.~Beck, A.~J\"ackle, G.~A.~Worth and H.-D.~Meyer, \textit{The multiconfiguration time-dependent Hartree (MCTDH) method: a highly efficient algorithm for propagating wavepackets}, Phys.\ Rep.\ \textbf{324}, 1 (2000).
    \item J.~Zanghellini, M.~Kitzler, T.~Brabec and A.~Scrinzi, \textit{Testing the multi-configuration time-dependent Hartree--Fock method}, J.\ Phys.\ B \textbf{37}, 763 (2004); J.~Caillat \textit{et al.}, \textit{Correlated multielectron systems in strong laser fields: A multiconfiguration time-dependent Hartree--Fock approach}, Phys.\ Rev.\ A \textbf{71}, 012712 (2005).
    \item A.~U.~J.~Lode, C.~L\'ev\^eque, L.~B.~Madsen, A.~I.~Streltsov and O.~E.~Alon, \textit{Colloquium: Multiconfigurational time-dependent Hartree approaches for indistinguishable particles}, Rev.\ Mod.\ Phys.\ \textbf{92}, 011001 (2020).
    \item O.~Leinonen, \textit{MCTDHF for simulations of multi-qubit systems}, in preparation; O.~Leinonen, PhD thesis, University of Oslo.
    \item R.~S.~Sutton and A.~G.~Barto, \textit{Reinforcement Learning: An Introduction}, 2nd ed., MIT Press (2018).
    \item J.~Schulman, F.~Wolski, P.~Dhariwal, A.~Radford and O.~Klimov, \textit{Proximal Policy Optimization Algorithms}, arXiv:1707.06347 (2017).
    \item M.~Bukov and F.~Marquardt, \textit{Reinforcement Learning for Quantum Technology}, arXiv:2601.18953 (2026).
    \item H.~Xu, J.~Li, L.~Liu, Y.~Wang, H.~Yuan and X.~Wang, \textit{Generalizable control for quantum parameter estimation through reinforcement learning}, npj Quantum Inf.\ \textbf{5}, 82 (2019).
    \item L.~Amico, R.~Fazio, A.~Osterloh and V.~Vedral, \textit{Entanglement in many-body systems}, Rev.\ Mod.\ Phys.\ \textbf{80}, 517 (2008).
    \item G.~Vidal and R.~F.~Werner, \textit{Computable measure of entanglement}, Phys.\ Rev.\ A \textbf{65}, 032314 (2002).
    \item W.~K.~Wootters, \textit{Entanglement of Formation of an Arbitrary State of Two Qubits}, Phys.\ Rev.\ Lett.\ \textbf{80}, 2245 (1998).
    \item V.~Coffman, J.~Kundu and W.~K.~Wootters, \textit{Distributed entanglement}, Phys.\ Rev.\ A \textbf{61}, 052306 (2000).
    \item D.~A.~Meyer and N.~R.~Wallach, \textit{Global entanglement in multiparticle systems}, J.\ Math.\ Phys.\ \textbf{43}, 4273 (2002).
    \item Z.-H.~Ma, Z.-H.~Chen, J.-L.~Chen, C.~Spengler, A.~Gabriel and M.~Huber, \textit{Measure of genuine multipartite entanglement with computable lower bounds}, Phys.\ Rev.\ A \textbf{83}, 062325 (2011).
    \item T.-C.~Wei and P.~M.~Goldbart, \textit{Geometric measure of entanglement and applications to bipartite and multipartite quantum states}, Phys.\ Rev.\ A \textbf{68}, 042307 (2003).
    \item L.~Pezz\`e, A.~Smerzi, M.~K.~Oberthaler, R.~Schmied and P.~Treutlein, \textit{Quantum metrology with nonclassical states of atomic ensembles}, Rev.\ Mod.\ Phys.\ \textbf{90}, 035005 (2018).
    \item F.~Belliardo, F.~Zoratti, F.~Marquardt and V.~Giovannetti, \textit{Model-aware reinforcement learning for high-performance Bayesian experimental design in quantum metrology}, Quantum \textbf{8}, 1555 (2024).
    \item M.~Bukov \textit{et al.}, \textit{Reinforcement Learning in Different Phases of Quantum Control}, Phys.\ Rev.\ X \textbf{8}, 031086 (2018); M.~Y.~Niu, S.~Boixo, V.~N.~Smelyanskiy and H.~Neven, \textit{Universal quantum control through deep reinforcement learning}, npj Quantum Inf.\ \textbf{5}, 33 (2019).
    \item F.~Benatti, R.~Floreanini, F.~Franchini and U.~Marzolino, \textit{Entanglement in indistinguishable particle systems}, Phys.\ Rep.\ \textbf{878}, 1 (2020); J.~Schliemann, J.~I.~Cirac, M.~Ku\'s, M.~Lewenstein and D.~Loss, \textit{Quantum correlations in two-fermion systems}, Phys.\ Rev.\ A \textbf{64}, 022303 (2001).
    \item M.~E.~Perruzza \textit{et al.}, \textit{Electrons on Helium and Entangled Quantum Sensors for Particle Physics}, in preparation (2026).
\end{enumerate}

\end{document}
